\section{Conversation chains}

Past students all say it: the exam is a conversation. The examiner takes a word from your answer
and asks about it. So every word you say is a door. This chapter turns that into a tool: each
answer ends with a \emph{hook}, one sentence that names the next topic. You open the door you
want, to a room you have prepared.

\begin{summary}{How to use the chains}
\textsf{Rule 1.} Answer in 2--3 sentences, then say the hook. Stop. Let the examiner walk through
the door. \
\textsf{Rule 2.} Only say words you can explain. ``Batch norm'' in your answer means a batch-norm
question in two seconds. \
\textsf{Rule 3.} If the examiner goes elsewhere, fine: jump to the matching step of another chain. \
\textsf{Rule 4.} Learn a chain as a \emph{story}, not as a list: each step is the answer to the
problem of the step before. \
Symbols: \textsf{\color{qcol}Q} question \textbullet{} \textsf{\color{qcol}A} your short answer
\textbullet{} \textsf{\color{piccol}$\to$ hook} the sentence that opens the next door.
\end{summary}

\newcommand{\cq}[1]{\par\smallskip\textsf{\color{qcol}Q.}\ #1}
\newcommand{\ca}[1]{\par\textsf{\color{qcol}A.}\ #1}
\newcommand{\hook}[1]{\par\textsf{\color{piccol}$\to$ hook:}\ \emph{#1}}

\newtcolorbox{chain}[1]{enhanced,breakable,colback=white,colframe=qcol!55,boxrule=0.5pt,arc=2pt,
  left=7pt,right=7pt,top=4pt,bottom=5pt,title={\sffamily #1},fonttitle=\small,coltitle=white,colbacktitle=qcol!85}

\begin{chain}{Chain 1 -- From data to batch normalization (Klambauer, Hochreiter)}
\pic{A story in one line: \emph{we want small risk} $\to$ \emph{we need gradients} $\to$
\emph{gradients vanish} $\to$ \emph{we fix the signal size} $\to$ \emph{with activations, with
initialization, with normalization}.}
\cq{What is the goal of learning?}
\ca{Small error on new data, the generalization error. We cannot measure it, so we minimize the
empirical risk $\Remp=\frac1N\sum_nL(y_n,g(\mathbf{x}_n;\mathbf{w}))$ on the training data.}
\hook{To minimize it we need a differentiable loss, so we use cross-entropy instead of the 0--1 loss.}

\cq{Why cross-entropy?}
\ca{It is differentiable, and minimizing it is maximum likelihood for a categorical output. With
softmax at the output the delta becomes simply $\hat{\mathbf{y}}-\mathbf{y}$.}
\hook{We then minimize it with gradient descent, in practice with mini-batches.}

\cq{Gradient descent and mini-batches?}
\ca{$\mathbf{w}\leftarrow\mathbf{w}-\eta\nabla\Remp$. A mini-batch gives a good, cheap gradient
estimate, runs in parallel on a GPU, and its noise helps to find flat minima.}
\hook{The gradient itself we get with backpropagation.}

\cq{Backpropagation?}
\ca{Forward pass, store $s$ and $a$. Delta at the output, then down:
$\delta_j=f'(s_j)\sum_i\delta_iw_{ij}$. Gradient $\partial L/\partial w_{ij}=\delta_ia_j$.}
\hook{Because the delta is multiplied by $f'$ and $W$ in every layer, it can vanish.}

\cq{Vanishing gradients?}
\ca{A product of many Jacobians shrinks or grows exponentially. With sigmoid $|f'|\le0.25$.
Found by Hochreiter in 1991.}
\hook{Three fixes: the activation function, the initialization, and normalization.}

\cq{Which activation helps?}
\ca{ReLU does not saturate for positive inputs. SELU, from Linz, even keeps mean 0 and variance 1
by itself: self-normalizing, with a stable fixed point.}
\hook{The same goal, a constant variance, is also the idea of good initialization.}

\cq{Initialization?}
\ca{Random, to break symmetry, with variance $1/n_{\mathrm{in}}$ (or $2/n_{\mathrm{in}}$ for ReLU)
so the signal keeps its size forward and backward.}
\hook{During training the statistics drift again, and batch normalization corrects that.}

\cq{Batch normalization?}
\ca{Standardize each pre-activation over the mini-batch,
$\hat s=(s-\mu_B)/\sqrt{\sigma_B^2+\epsilon}$, then learn scale and shift, $\gamma\hat s+\beta$.
Training is faster and more stable, and it regularizes a little. Layer norm does the same over the
features instead of the batch, which suits RNNs and Transformers.}
\hook{SELU reaches the same effect without any batch statistics, which is why it was developed here.}
\end{chain}

\begin{chain}{Chain 2 -- From RNN to xLSTM (Hochreiter)}
\pic{\emph{RNNs forget} $\to$ \emph{LSTM remembers on a conveyor belt} $\to$ \emph{Transformers
look at everything, but pay $N^2$} $\to$ \emph{xLSTM brings memory back at linear cost}.}
\cq{What is an RNN?}
\ca{A network with a hidden state that is fed back: $\mathbf{s}(t)=W\mathbf{x}(t)+R\mathbf{a}(t-1)$.
Same weights at every step, any sequence length, Turing complete.}
\hook{We train it with backpropagation through time.}

\cq{BPTT?}
\ca{Unroll in time and backprop: $\bm\delta(t)=\mathrm{diag}(f')(V^\top\partial L(t)/\partial\hat{\mathbf{y}}+R^\top\bm\delta(t+1))$.
Memory grows with the sequence length; RTRL avoids that but costs $O(N^4)$ per step.}
\hook{The problem is that the error from far in the past vanishes.}

\cq{Why does it vanish?}
\ca{At every step it is multiplied by $R^\top\mathrm{diag}(f')$; after $k$ steps by the $k$-th power.
Clipping helps against exploding, not vanishing.}
\hook{LSTM solves it with the constant error carousel.}

\cq{The constant error carousel?}
\ca{A memory cell updated by addition, $\mathbf{c}(t)=\mathbf{f}\odot\mathbf{c}(t-1)+\mathbf{i}\odot\mathbf{z}$,
so $\partial\mathbf{c}(t)/\partial\mathbf{c}(t-1)=\mathrm{diag}(\mathbf{f})\approx I$. The error rides
the belt unchanged. Gates decide what to write, keep and read.}
\hook{The forget gate was added later, and GRU simplified the gates.}

\cq{Forget gate and GRU?}
\ca{Gers et al.\ 2000: the forget gate lets the cell reset on long streams. GRU couples input and
forget into one update gate and adds a reset gate; fewer parameters, similar results.}
\hook{LSTMs still read step by step; attention removed that bottleneck.}

\cq{Attention and Transformers?}
\ca{$\mathrm{softmax}(QK^\top/\sqrt{d_k})V$: every token looks at every token in one step, fully
parallel. The price is $O(N^2)$ in the sequence length.}
\hook{xLSTM tries to keep the parallel training but with linear cost.}

\cq{xLSTM?}
\ca{Exponential gating with a normalizer, so storage decisions can be revised; and mLSTM, a matrix
memory $C_t=f_tC_{t-1}+i_t\mathbf{v}_t\mathbf{k}_t^\top$, parallel in training, linear in inference.}
\hook{And the attention formula itself is the update rule of a modern Hopfield network.}
\end{chain}

\begin{chain}{Chain 3 -- Generalization and regularization (Klambauer)}
\pic{\emph{Fit the data} $\to$ \emph{but not too well} $\to$ \emph{measure it honestly}
$\to$ \emph{control it}.}
\cq{What is overfitting?}
\ca{Training error low, test error high: the model learned the noise. Too much capacity for the data.}
\hook{Capacity can be measured, for example by the VC dimension.}
\cq{VC dimension?}
\ca{The largest point set the model class can shatter; lines in 2D: 3. Bounds say the gap between
test and training error grows with it and shrinks with $N$.}
\hook{In practice we do not use the bound, we estimate the test error directly.}
\cq{How?}
\ca{Held-out test set, or $L$-fold cross-validation, which is almost unbiased and uses all data.}
\hook{If the estimate shows overfitting, we regularize.}
\cq{Regularization?}
\ca{Weight decay $\frac\lambda2\lVert\mathbf{w}\rVert^2$, dropout, early stopping, data augmentation,
Flat Minimum Search. They all limit effective capacity.}
\hook{Interestingly, very large networks can generalize even without it: double descent.}
\cq{Double descent?}
\ca{Past the interpolation point the test error falls again as the model grows; the classical
U-curve is only the first half.}
\end{chain}

\begin{chain}{Chain 4 -- From your thesis into ML (Klambauer, chair)}
\pic{\emph{rCFD is data reuse} $\to$ \emph{data reuse means generalization} $\to$ \emph{generalization
fails outside the training regime} $\to$ \emph{a learned model would face the same limit}.}
\cq{Why is your thesis AI?}
\ca{The expensive solver produces data once, and a cheap model reuses it. Choosing the next frame
is a nearest-neighbour search over recorded states.}
\hook{So the key question is the same as in ML: how far does it generalize?}
\cq{How far does it generalize?}
\ca{Within its regime. The adiabatic recording fails under the cooled lid because cold plumes
create a motion it never saw: out-of-distribution.}
\hook{A learned surrogate would have the same limit, but it could learn the missing correction.}
\cq{Which model would you use?}
\ca{A graph neural network on the mesh: cells are nodes, faces are edges, message passing like the
face exchange. Loss: distance to the reference plus an energy-balance penalty.}
\hook{The penalty matters, because conservation is what rCFD guarantees by construction.}
\cq{Why conservation?}
\ca{My corrections are conservative: every face exchange appears twice with opposite signs, so the
heat in the tank only moves. A learned model without that constraint can drift in energy.}
\end{chain}

\begin{chain}{Chain 5 -- From numerical diffusion to your limiter (Pirker)}
\pic{\emph{Schemes smear} $\to$ \emph{explicit steps have limits} $\to$ \emph{my cap is such a limit}
$\to$ \emph{un-smearing needs a limiter}.}
\cq{What is numerical diffusion?}
\ca{First-order upwind is stable but adds false diffusion that smears fronts, error $\propto\Delta x$.
Central differencing is sharper but oscillates when the grid P\'eclet number exceeds 2.}
\hook{In my replay the smearing comes from the face exchange, which is an explicit step.}
\cq{Limits of explicit steps?}
\ca{Convection needs $Co=u\Delta t/\Delta x<1$, diffusion $d=\Gamma\Delta t/\Delta x^2<0.5$.}
\hook{This is exactly why the recorded diffusivity could not act.}
\cq{Why not?}
\ca{The exchange is capped at $\mathrm{Di}\le0.5$; the recorded field asked for about 4. Every face hit
the cap, so it behaved like the constant $\fsd=0.5$.}
\hook{To undo the smearing I needed anti-diffusion, and that needs a limiter.}
\cq{Your limiter?}
\ca{Negative diffusion is ill-posed and blows up at the grid scale. So I move heat up the gradient
only as far as both cells' local bounds allow, $\Delta q_f=s\min(Q_h^+,Q_c^-)$, like flux-corrected
transport: conservative, no new extremum.}
\hook{And I weight it towards the vertical, because the front is horizontal.}
\end{chain}

\subsection*{Doors to open and doors to keep closed}

\begin{center}\small
\begin{tabularx}{\linewidth}{@{}XX@{}}
\toprule
Say it (you have a room behind it) & Avoid unless asked (deep rooms) \\ \midrule
empirical risk, cross-entropy, mini-batch & natural gradient, Fisher information \\
delta error, vanishing gradient, SELU & neural tangent kernel \\
constant error carousel, forget gate, GRU & multidimensional LSTM details \\
attention, $\sqrt{d_k}$, positional encoding & rotary embeddings in depth \\
flux-corrected transport, limiter, conservative & Harten's TVD theory \\
out-of-distribution, regime, nearest neighbour & Koopman operators, propagators in depth \\
\bottomrule
\end{tabularx}
\end{center}
